\section{Osteoporosis Screening --- DEXA Scan}


\subsubsection*{What Is This Test?}
This test measures specific substances in your blood, urine, or other body fluids to help diagnose, monitor, or screen for medical conditions. Your doctor will explain why this particular test was chosen for your situation.

\subsubsection*{Alternative Names}
DEXA scan, DXA, bone density test, dual-energy X-ray absorptiometry, bone densitometry, central DXA, peripheral DXA


\subsubsection*{Principle}
DXA uses two X-ray beams at different energies (typically 70 and 140 kVp). The differential attenuation of the two energies distinguishes bone from soft tissue. Bone mineral density (BMD) in g/cm2 is calculated. The T-score compares the patient's BMD to a young adult reference population.

\subsubsection*{History}
Bone densitometry evolved from plain X-ray techniques in the 1960s. Dual-energy X-ray absorptiometry (DXA) was developed in the 1980s and became the gold standard for osteoporosis diagnosis after the WHO defined T-score criteria in 1994.

\subsubsection*{Why This Test Is Ordered}
Your bones get thinner as you age, especially after menopause. A DEXA scan measures bone density to find out if you are at risk for fractures before one actually happens. It answers the question: ``Are your bones strong enough to survive a fall?''

\subsubsection*{Is This a Primary or Secondary Test}
Primary screening test for osteoporosis.

\subsubsection*{How DXA performance is assessed}
DXA provides a bone-density measurement when positioning and quality control are adequate. Its clinical value is assessed through fracture-risk prediction and change beyond measurement error, rather than sensitivity for one disease.

\subsubsection*{How This Test Is Performed}
You lie on a padded table while a machine scans your hip and spine with a low-dose X-ray. The scan takes about 15 minutes. There is no needle, no injection, and no pain. You stay fully clothed.

\subsubsection*{What Preparation Is Needed}
Some facilities ask patients to avoid calcium supplements for a period before DXA; follow the local appointment instructions rather than stopping a prescribed supplement on your own. Wear clothing without metal fasteners when requested, and tell the imaging team if you are pregnant or may be pregnant.

\subsubsection*{Contraindications and Precautions}
Pregnancy is not an automatic contraindication; the examination should be justified and optimized, and alternatives considered when appropriate. For contrast-enhanced studies, review the indication, prior severe reactions, and renal function according to local protocol.

\subsubsection*{Risks}
Risks associated with X-ray imaging include:
\begin{itemize}
    \item \textbf{Radiation exposure}: Exposure to low levels of ionizing radiation, with cumulative lifetime dose optimization guided by the ALARA (As Low As Reasonably Achievable) principle.
    \item \textbf{Contrast reaction}: Allergic reactions can occur, including rare severe reactions. Kidney injury after iodinated contrast is uncommon and risk depends on baseline kidney function, illness, and the clinical context.
\end{itemize}


You can resume normal activities immediately. Results usually take a few days. Your doctor will compare your bone density to two numbers: the peak bone density of a healthy 30-year-old (your T-score) and the average for someone your age and sex (your Z-score).

\subsubsection*{Understanding Your Results}

\subsubsection*{Normal Results}
Normal: T-score $\ge$ -1.0. Osteopenia: T-score -1.0 to -2.5. Osteoporosis: T-score $\le$ -2.5.

\begin{center}
\begin{tabular}{ll}
\toprule
\textbf{Population} & \textbf{Reference Range} \\
\midrule
Adult Female & T-score $\ge$ -1.0 \\
Adult Male & T-score $\ge$ -1.0 \\
Children & Not applicable \\
Elderly & T-score $\ge$ -1.0 (age-adjusted Z-score) \\
Pregnancy & Not applicable \\
\bottomrule
\end{tabular}
\end{center}
\subsubsection*{Follow-up Tests}
When a test result is not within the normal range, your doctor may recommend additional testing to confirm the finding, identify the underlying cause, and guide treatment decisions:
\begin{itemize}
    \item \textbf{Serum calcium, phosphate, and PTH:} Measured simultaneously to classify hypercalcemia (PTH-dependent vs. PTH-independent) and guide further workup.
    \item \textbf{25-hydroxyvitamin D and 1,25-dihydroxyvitamin D:} To differentiate vitamin D deficiency from vitamin D resistance or renal vitamin D metabolism disorders.
    \item \textbf{24-hour urinary calcium:} To evaluate renal calcium handling and differentiate absorptive from renal hypercalciuria.
    \item \textbf{Bone mineral density scan (DEXA):} To assess bone density and fracture risk when chronic PTH or vitamin D abnormalities are present.
\end{itemize}
\clearpage
